% Generated by figures/13_14_24_25_granger_tables.R (score: a2) -- do not edit by hand.
% Requires \usepackage{booktabs} and \usepackage{threeparttable}.
\begin{table}[htbp]
\centering
\begin{threeparttable}
\caption{Granger causality between US and Chinese sentiment, by topic}
\label{tab:granger-topics}
\small
\begin{tabular}{l rrr @{\hspace{1.5em}} rrr}
\toprule
 & \multicolumn{3}{c}{All years} & \multicolumn{3}{c}{After 2015} \\
\cmidrule(lr){2-4}\cmidrule(lr){5-7}
Topic & US as & China as & Months & US as & China as & Months \\
 & response & response &  & response & response &  \\
\midrule
Domestic periphery \& human rights$^{b}$ & 0.268 & 0.066 & 87 & 0.630 & 0.785 & 42 \\
Multilateral negotiations & 0.443 & 0.243 & 201 & \textbf{0.043} & 0.139 & 80 \\
Trade \& finance & \textbf{0.043} & 0.555 & 217 & \textbf{0.013} & 0.898 & 79 \\
Taiwan$^{b}$ & 0.426 & 0.636 & 69 & 0.963 & \textbf{0.041} & 22 \\
South China Sea$^{a}$ & 0.376 & 0.938 & 45 & 0.441 & 0.975 & 36 \\
Global health \& epidemics$^{a}$ & 0.807 & \textbf{0.042} & 27 & 0.145 & 0.504 & 22 \\
Russia & 0.379 & 0.352 & 147 & \textbf{0.045} & \textbf{0.027} & 69 \\
\bottomrule
\end{tabular}
\begin{tablenotes}[flushleft]
\footnotesize
\item \textit{Notes:} Wald test $p$-values of bivariate Granger causality tests with 3 lags of each series. ``US as response'' tests whether lagged Chinese sentiment improves the prediction of US sentiment beyond lagged US sentiment; ``China as response'' tests the reverse. Bold: $p < 0.05$. Sentiment is the mean hybrid score (approach 2) of the documents in each period. Each pair of series is first tested for a unit root (augmented Dickey--Fuller, constant and trend, 5\% level, full sample). If both series are stationary, the test is run in levels; if neither is, on first differences; if only one is, by the Toda--Yamamoto procedure (both series in levels with one additional, untested lag of each). All $p$-values are from $F$-tests. Each series is the monthly mean score of the documents matching a topic's keywords (full US and Chinese data). Months in which either country has no document on the topic are dropped before lagging; ``Months'' counts the months left. The right panel uses only months after 1 January 2015. The test for each topic is chosen on all months and kept for the right panel. $^{a}$ First differences. $^{b}$ Toda--Yamamoto. Other rows: levels.
\end{tablenotes}
\end{threeparttable}
\end{table}
